A self-adaptive system rarely faces one fixed environment: layouts, sizes and mission lengths change, and
synthesising a new controller for every variant at run time is costly. We ask whether a large language model (LLM)
can write a few readable controllers for a whole family of Markov decision processes (MDPs) that a probabilistic
model checker then certifies on each instance the system meets, including instances larger than any seen in
training. The LLM writes ordered rule lists over a domain vocabulary of features and instance constants; a loop
verifies them with PRISM on a few small training instances and feeds the results back, and the verifier refines the
lists by resolving the choices they leave open into the action order that does best on the training instances.
At run time the system certifies its stored lists on the instance at hand and uses one that passes. On random
gridworlds, on three parameterised gridworld families and on the Pac-Man benchmark, the four lists of a single
training run cover \todo{all} held-out validation grids, and three runs cover \todo{49 of 60} family members, against
\todo{37} for the best baseline: decision trees learned from optimal strategies (as in 1--2--3--Go!) or PPO, both
given the same features. Prompting variants change little; refinement by the verifier is what makes the lists
transfer, at about one cent per training set with Claude Haiku~5.5.
